\clearpage
\fancyhead[L]{\small HERSCHEL AND SOUTH. OBSERVATIONS OF DOUBLE STARS.}
\hypertarget{native-647-lower}{}
\begin{center}\small Göttingische gelehrte Anzeigen. Stück 201. Seiten 2001..2004. 18. Dezember 1826.\\[3pt]\rule{.72\linewidth}{.3pt}\end{center}

\emph{Observations of the apparent distances and positions of \textup{380} double and triple stars made in the years \textup{1821}, \textup{1822} and \textup{1823}, and compared with those of other astronomers.} London 1825. By \textsc{J. F. W. Herschel} and \textsc{James South}. 412 Seiten in 4., nebst 4 Kupfertafeln.

\emph{Observations of the apparent distances and positions of \textup{458} double and triple stars made in the years \textup{1823}, \textup{1824} and \textup{1825}.} Ebenda 1826. By \textsc{James South}. 391 und XVIII Seiten in 4.

\begingroup\baselineskip=12.5pt\relax
Die Geschichte der Astronomie zeigt uns oft, dass wichtige unerwartete Aufschlüsse aus Beobachtungen gewonnen wurden, die lange zuvor mit mühevollem Fleiß angestellt und aufgezeichnet worden waren, ohne einen solchen Erfolg zu ahnen. Die Doppelsterne sind ein bemerkenswertes Beispiel dafür. Vor \textsc{Herschel} kannte man zwar schon eine, wenn auch vergleichsweise kleine, Anzahl von Doppelsternen, schenkte ihnen jedoch keine besondere Aufmerksamkeit. Man ging von der \SourceLeaf{648}Meinung aus, dass das Phänomen eines Doppelsterns nur das Spiel eines Zufalls sei, der zwei Fixsterne von unserem Sonnensystem aus gesehen beinahe in dieselbe Richtung gestellt habe, wobei ihre wirklichen Entfernungen im Allgemeinen sehr ungleich sein könnten. \textsc{Herschel} machte als Erster auf die Vorteile aufmerksam, die die fortgesetzte Beobachtung eines Doppelsterns für die Bestimmung der jährlichen Parallaxe bieten könnte, eine Methode, deren Brauchbarkeit offenbar ganz von jener Voraussetzung abhängt. \textsc{Herschel} zeichnete bei seiner Durchmusterung des Sternenhimmels alle Doppelsterne auf, die ihm begegneten, beobachtete ihre gegenseitigen Stellungen mit großer Genauigkeit und stellte so ein Verzeichnis von mehr als 700 Doppel- und Mehrfachsternen zusammen. Darunter sind freilich viele, denen man wegen der beträchtlichen Entfernung der einzelnen Sterne voneinander nur im weiteren Sinne den Namen Doppelsterne geben kann, obwohl es in der Natur der Sache liegt, dass sich hinsichtlich dieser Entfernung keine scharfen Grenzen ziehen lassen.

Zwanzig Jahre später unterzog \textsc{Herschel} die Doppelsterne einer erneuten Prüfung. Bei vielen von ihnen bemerkte er in der gegenseitigen Stellung der einzelnen Sterne solche Veränderungen, dass er zu der Überzeugung gelangte, die frühere Auffassung einer zufälligen Bildung der Doppelsterne sei nicht allgemein haltbar und wenigstens viele von ihnen könnten \emph{wirklich} nahe beieinanderstehende Sternpaare sein, die sich nach bestimmten, in den Naturkräften begründeten Gesetzen umeinander bewegen.

Diese höchst wichtige Folgerung gewann durch die sehr große Anzahl von Doppelsternen, die am Himmel beobachtet werden, erheblich an Wahrscheinlichkeit. Man kann aber getrost behaupten, dass es all dessen eigentlich kaum bedurft hätte, um die Schwäche der früheren Voraussetzung offenzulegen. Wenn wir in der Tat etwa nur den einen Doppelstern Castor betrachten, der aus zwei Sternen besteht, die sich in ihrer Helligkeit nicht sehr unterscheiden, beide ungefähr der dritten Größe angehören und etwa fünf Bogensekunden voneinander abstehen, dann finden wir die Wahrscheinlichkeit der Bildung eines solchen Doppelsterns unter der Voraussetzung einer ganz zufälligen Verteilung sämtlicher Fixsterne bis zur dritten Größe auf der Himmelskugel so außerordentlich klein, dass wir es im gewöhnlichen Leben bei ähnlichen Ereignissen für ungereimt halten, an einer \emph{nicht} zufälligen Entstehung zu zweifeln. Unmöglich ist es freilich nicht, dass unter der großen Menge der Doppelsterne einige von zufälliger Entstehung sein mögen. Doch muss man annehmen, dass die meisten Doppelsterne mit geringerem Abstand ihrer Einzelsterne tatsächlich zusammengehörige Systeme bilden, deren gegenseitige Bewegungen ein neues Feld interessanter Beobachtungen eröffnet haben. Diese Bewegungen sind aber so langsam, dass bei den meisten selbst in einem Menschenalter die Verschiebung nur gering ist. Doch haben sich auch einige Doppelsterne gefunden, bei denen sie verhältnismäßig schon in sehr kurzen Zeiträumen sichtbar wird. Auf jeden Fall ist klar, dass es von großer Wichtigkeit ist, die Doppelsterne von Zeit zu Zeit erneut zu untersuchen, um den künftigen Jahrhunderten das Material für eine sehr viel umfassendere Kenntnis des Weltgebäudes bereitzustellen.

Dies haben nun die Herren \textsc{Herschel} d. J., Sohn des berühmten Beobachters, und \textsc{South} in den vorliegenden Werken auf musterhafte Weise getan. Das erste Werk enthält die Beobachtungen, die diese beiden Astronomen gemeinsam mit einem fünffüßigen und einem siebenfüßigen Äquatoreal angestellt haben. Das zweite enthält diejenigen, die der letztere Astronom allein gemacht hat, teils mit beiden Instrumenten in England, teils mit dem siebenfüßigen Äquatoreal in Frankreich an einer vorübergehend in Passy errichteten Sternwarte. Wir haben hier also einen reichen Schatz an Beobachtungen über 838 Doppelsterne aus den Jahren 1821..1825 und zugleich ihre Zusammenstellung mit den früheren Beobachtungen des älteren \textsc{Herschel}, die zum Teil aus dem handschriftlichen Nachlass berichtigt und ergänzt wurden, sowie mit den nicht minder wertvollen Beobachtungen von \textsc{Struve}. Bei den meisten Doppelsternen sind die bisher beobachteten Veränderungen in den Stellungen zwar nur klein, wenn auch \SourceLeaf{649}eindeutig. Die Umlaufzeit einer Sonne um die andere wird im Allgemeinen nach Jahrhunderten gemessen. Indessen finden sich bei manchen schon so bedeutende Änderungen, dass wir dadurch eine ungefähre Vorstellung von der Länge der Periode erhalten. Beim bemerkenswerten Stern 61 Schwan zum Beispiel hat die Veränderung des Richtungswinkels vom Jahr 1753 bis 1825 schon 53 Grad betragen. Bei Castor waren es von 1759 bis zu diesem Jahr 63 Grad, bei $\sigma$ Krone seit 1781 90 Grad, bei $\xi$ im Großen Bären im selben Zeitraum 259 Grad und bei 70 Ophiuchus seit 1779 sogar 302 Grad. Der begrenzte Raum verbietet uns, noch manches andere Bemerkenswerte hervorzuheben. Dazu gehören Sterne, die früher doppelt erschienen und sich gegenwärtig vollkommen decken, dreifache Sterne, die Ein System bilden und sich gegenseitig bewegen, Sterne, die zwölf Bogenminuten voneinander abstehen und doch sehr wahrscheinlich nur Ein System bilden, und dergleichen. [S. Neue Aussicht zur Erweiterung u. s. f. \hyperlink{native-186}{oben S. 181.} \textsc{Schering}.]
\par\clearpage\endgroup
